\scanpage{31}
% archive sheet number: 30 (inverted)
\mspage{- d -}\hfill n°383

\begin{center}\uline{Commentaires}\end{center}

On a groupé, dans un \S~1 préliminaire, la plupart des lemmes géométriques qui venaient précédemment rompre la suite des idées. On a défini les chambres pour des familles localement finies d'hyperplans, ce qui est utile au \S~3 ; tout marche. Le c) du lemme 2 est inutile pour la suite, mais ça rassure.

Au \S~2, on a suivi Tits pas à pas. C'est assez pénible. On a ajouté de jolis compléments extraits de Tits, presque tous inutiles pour la suite.

Naturellement l'effort fait au \S~2 est récompensé au \S~3, où la théorie des groupes finis engendrés par des réflexions marche comme sur des roulettes. Cependant il y a, dans le cas affine, un vilain calcul pour démontrer le th.3 ; Bourbaki sera sûrement assez astucieux pour le court-circuiter.

Le \S~4 des systèmes de racines a été agrémenté des compléments dûs à Borel, ainsi que de quelques figures. On a le sentiment d'avoir notablement simplifié le n° des ensembles clos de racines.

Le rédacteur a profité de ce qu'il y avait de l'algèbre commutative au \S~5 pour plus chambouler les choses que ne lui prescrivait la Tribu. Il n'a rien redémontré de ce qui se trouve dans le Livre d'Algèbre Commutative, ni de ce qui s'y trouvera, ni de ce qui devrait s'y trouver. Il a remplacé un lemme bizarre de l'état précédent par un autre qui l'est moins. Ayant eu besoin du théorème de Maschke, on l'a démontré en lemme, et on s'est demandé par quel mystère insondable il n'est pas en Algèbre VIII. Pour démontrer que l'algèbre des invariants est une algèbre de polynômes, on s'est borné au cas où l'anneau de base est un corps d'exposant caractéristique étranger à ce qu'il faut ; des notes prises en Congrès semblent dire que ça marche aussi sur un anneau convenable, mais le rédacteur s'est d'abord échiné en vain, puis dégonflé.

Le \S~6 des transformations de Coxeter a été fort peu modifié.

Au \S~7 (Classification), on a utilisé le Séminaire Sophus Lie pour simplifier la recherche des schémas admissibles au moyen de lemmes bien choisis. Pour la construction des systèmes de types $A_r,B_r,C_r,D_r$ et $G_2$, on a froidement copié Dixmier (cependant, en remplaçant par nu raisonnement un calcul relatif à $D_r$, on a donné raison à Dixmier contre Dynkin). Une astuce, vendue par Serre, a permis d'éviter les vérifications assommantes relatives à $E_{8},E_{7}$ et $E_{6}$ ; une astuce analogue a été utilisée pour $F_{4}$.
